\documentclass[../../main.tex]{subfiles}
\begin{document}

{\bf [解]}

パラメータ$\theta\, (0\le\theta<2\pi)$を用いて
$C$上の点$P$は $P(\sqrt{\alpha}\cos\theta, \sqrt{\beta}(\sin\theta+1))$ とかける．また題意より$\alpha,\beta >0$である．

\begin{figure}[htb]
  \centering
  \begin{tikzpicture}[scale=3, >=stealth, every node/.style={font=\small}]

    % --- 1. 定数・幾何パラメータ ---
    % D: 単位円盤 x^2 + y^2 <= 1
    % C: 中心 (0, sqrt(beta)), 長半径・短半径 sqrt(alpha), sqrt(beta)
    % 例として (2) で最大面積を与えるパラメータ付近を採用:
    % alpha = 2/3 ≈ 0.667 => sqrt(alpha) ≈ 0.816
    % beta = 2/9 ≈ 0.222  => sqrt(beta)  ≈ 0.471
    % 中心 y_0 = sqrt(beta) ≈ 0.471
    \def\aParam{0.6667}
    \def\bParam{0.2222}
    \pgfmathsetmacro{\radA}{sqrt(\aParam)} % ≈ 0.816
    \pgfmathsetmacro{\radB}{sqrt(\bParam)} % ≈ 0.471
    \pgfmathsetmacro{\yC}{\radB}           % 中心 y 座標 ≈ 0.471

    % 代表点 P (theta = 50°)
    \def\angTheta{50}
    \pgfmathsetmacro{\xP}{\radA * cos(\angTheta)}
    \pgfmathsetmacro{\yP}{\radB * (sin(\angTheta) + 1.0)}

    \coordinate (O)     at (0, 0);
    \coordinate (Center) at (0, \yC);
    \coordinate (P)     at (\xP, \yP);

    % --- 2. 領域 D (単位円盤) の塗りつぶしと輪郭線 ---
    \fill[gray!10, opacity=0.7] (O) circle (1.0);
    \draw[thick, gray!70] (O) circle (1.0);
    \node[gray!80!black, font=\footnotesize] at (-0.68, -0.68) {領域 $D: x^2+y^2 \le 1$};

    % --- 3. 楕円 C (中心 (0, sqrt(beta))) の塗りつぶしと輪郭線 ---
    % 楕円の方程式: x^2/alpha + (y - sqrt(beta))^2/beta = 1
    \fill[blue!15, opacity=0.7] (Center) ellipse [x radius=\radA, y radius=\radB];
    \draw[very thick, blue!85!black] (Center) ellipse [x radius=\radA, y radius=\radB];
    \node[blue!90!black, font=\footnotesize\bfseries] at (-0.28, 0.47) {$C$};

    % --- 4. 座標軸 ---
    \draw[->, thick, gray!80!black] (-1.35, 0) -- (1.45, 0) node[right] {$x$};
    \draw[->, thick, gray!80!black] (0, -1.25) -- (0, 1.4) node[above] {$y$};
    \node[below left=2pt] at (O) {$O$};

    % 単位円の目盛り (±1)
    \node[below=2pt, font=\scriptsize] at (1.0, 0) {$1$};
    \node[below=2pt, font=\scriptsize] at (-1.0, 0) {$-1$};
    \node[left=2pt, font=\scriptsize]  at (0, 1.0) {$1$};
    \node[left=2pt, font=\scriptsize]  at (0, -1.0) {$-1$};
    \fill (1.0, 0)  circle (0.5pt);
    \fill (-1.0, 0) circle (0.5pt);
    \fill (0, 1.0)  circle (0.5pt);
    \fill (0, -1.0) circle (0.5pt);

    % --- 5. 楕円の中心と軸半径の寸法表示 ---
    % 中心 (0, sqrt(beta))
    \fill[teal!85!black] (Center) circle (1.0pt);
    \node[left=2pt, font=\scriptsize, teal!90!black] at (Center) {$(0, \sqrt{\beta})$};

    % 半径 sqrt(alpha) (水平方向)
    \draw[<->, thick, teal!85!black] (Center) -- (\radA, \yC);
    \node[teal!90!black, font=\scriptsize, above] at ($ (Center)!0.5!(\radA, \yC) $) {$\sqrt{\alpha}$};

    % 半径 sqrt(beta) (垂直方向)
    \draw[<->, thick, teal!85!black] (Center) -- (0, {\yC + \radB});
    \node[teal!90!black, font=\scriptsize, right] at ($ (Center)!0.5!(0, {\yC + \radB}) $) {$\sqrt{\beta}$};

    % --- 6. 動点 P と原点からの距離ベクトル OP ---
    \draw[->, line width=1.4pt, red!85!black] (O) -- (P);
    \fill[red!85!black] (P) circle (1.2pt);
    \node[above right=1pt, red!90!black, font=\footnotesize]
    at (P) {$P\left(\sqrt{\alpha}\cos\theta, \sqrt{\beta}(\sin\theta+1)\right)$};

    % 距離 |OP| <= 1 の注記
    \node[red!90!black, font=\scriptsize, left=1pt] at ($ (O)!0.55!(P) $) {$|OP| \le 1$};

  \end{tikzpicture}
  \caption{単位円盤 $D$ とそれに含まれる楕円 $C$ および動点 $P$}
  \label{fig:ellipse_in_disk_overview}
\end{figure}

\paragraph{(1)}

$C$が$D$に含まれる条件は，$C$上の任意の点$P$が原点との距離が$1$以下であること，
すなわち
\begin{align}
  (\sqrt{\alpha}\cos\theta)^2 + (\sqrt{\beta}(\sin\theta+1))^2 &\le 1
\end{align}
が任意の$\theta$で成立することである．
展開して整理して
\begin{align}
  \alpha (1 - \sin^2\theta) + \beta (\sin\theta+1)^2 - 1 &\le 0 \\
  \therefore
  (\beta - \alpha) \sin^2\theta + 2\beta \sin\theta + \alpha + \beta - 1 &\le 0 \label{eq:1}
\end{align}
である．\cref{eq:1}の左辺を$f(\sin\theta)\, (-1\le \sin\theta \le 1)$とおくと
\begin{align}
  f(x) = (\beta - \alpha) x^2 + 2\beta x + \alpha + \beta - 1 \label{eq:2}
\end{align}
だから，一階微分は
\begin{align}
  f'(x) = 2(\beta - \alpha)x + 2\beta \label{eq:3}
\end{align}
である．$f'(x)$はたかだか一次関数で，両端での値は
\begin{align}
  \begin{cases}
    f'(1) = 2(2\beta-\alpha) \\
    f'(-1) = 2\alpha >0
  \end{cases}
\end{align}
である．$f'(-1)>0$だから，$f'(1)$の正負に応じて$f(x)$の増減表は変化する．

\subparagraph{$1^\circ \ f'(1) \ge 0 \iff 2\beta \ge \alpha$ の時}

$f(x)$は高々1次関数であるから，$[-1, 1]$で$f'(x) \ge 0$であり，$f(x)$は単調増加であり，増減表は以下のようになる．

\begin{table}[htb]
  \centering
  \caption{$f(x)$ の増減表 ($2\beta \ge \alpha$ のとき)}
  \renewcommand{\arraystretch}{1.4}
  \begin{tabular}{|c|c|c|c|}
    \hline
    $x$    & $-1$       & $\cdots$   & $1$       \\ \hline
    $f'$   & $+$        & $+$        & $+$       \\ \hline
    $f$    & $f(-1)$    & $\nearrow$ & $f(1)$    \\ \hline
  \end{tabular}
\end{table}

したがって，この区間での$f(x)$の最大値は$f(1)$だから，$f(1) \le 0$ ならば\cref{eq:1}をみたす．よって求める条件は
\begin{align}
  f(1) = 4\beta - 1 \le 0 \iff \beta \le \dfrac{1}{4} \label{eq:4}
\end{align}
である．

\subparagraph{$2^\circ \ f'(1) \le 0 \iff 2\beta \le \alpha$ の時}

$f'(x)$は高々$S$の1次関数だから$[-1,1]$で一回だけ符号を変える．
$\beta \ne \alpha$ だから，$f(x)$の増減表は以下のようになる．

\begin{table}[htb]
  \centering
  \caption{$f(x)$の増減表}
  \begin{tabular}{|c|c|c|c|c|c|}
    \hline
    $x$ & $-1$ & $\cdots$ & $-\dfrac{\beta}{\beta-\alpha}$ & $\cdots$ & $1$ \\
    \hline
    $f'$ & $+$ & $+$ & $0$ & $-$ & $-$ \\
    \hline
    $f$ & & $\nearrow$ & & $\searrow$ & \\
    \hline
  \end{tabular}
\end{table}

したがって，$f\left(-\dfrac{\beta}{\alpha-\beta}\right) \le 0$ なら条件をみたす．
\begin{align}
  f\left(-\dfrac{\beta}{\alpha-\beta}\right) = \dfrac{\beta^2}{\alpha-\beta} + \alpha + \beta - 1 &\le 0
\end{align}
$0 < 2\beta \le \alpha$ より $\alpha - \beta > 0$ だから，両辺 $\alpha - \beta$ をかけて，求める条件は
\begin{align}
  \beta \le -\alpha^2 + \alpha \label{eq:5}
\end{align}
である．
\casesend

\cref{eq:4,eq:5}をまとめて求める条件は，
\begin{align}
  \begin{cases}
    \alpha \le 2\beta \text{ の時， } \beta \le \dfrac{1}{4} \\
    2\beta \le \alpha \text{ の時， } \beta \le -\alpha^2 + \alpha
  \end{cases} \label{eq:6}
\end{align}
である．これを図示して下図斜線部（境界は$\alpha, \beta$軸のみ含まず）である．

\begin{figure}[htb]
  \centering
  \begin{tikzpicture}[scale=6, >=stealth, every node/.style={font=\small}]

    % --- 1. 定数・座標パラメータ ---
    % 接続点: alpha = 1/2, beta = 1/4 (直線 beta = 1/2*alpha と放物線 beta = -alpha^2 + alpha の交点)
    \coordinate (O)       at (0, 0);
    \coordinate (P_joint) at (0.5, 0.25);
    \coordinate (P_alpha1) at (1.0, 0);
    \coordinate (P_beta14) at (0, 0.25);

    % --- 2. 求める領域のハッチング塗りつぶし ---
    % 0 < alpha <= 1/2 では beta <= 1/4
    % 1/2 <= alpha < 1 では beta <= -alpha^2 + alpha
    \fill[pattern=north east lines, pattern color=blue!65, opacity=0.75]
    (0, 0) -- (0, 0.25) -- (P_joint)
    -- plot[domain=0.5:1.0, samples=50] (\x, {-\x*\x + \x})
    -- cycle;

    % --- 3. 領域の境界線 ---
    % (i) beta = 1/4 (0 <= alpha <= 1/2) (実線)
    \draw[very thick, blue!85!black] (0, 0.25) -- (P_joint);

    % (ii) 放物線 beta = -alpha^2 + alpha (1/2 <= alpha <= 1) (実線)
    \draw[very thick, blue!85!black, domain=0.5:1.0, samples=60]
    plot (\x, {-\x*\x + \x});

    % 放物線の延長点線 (0 <= alpha <= 0.5)
    \draw[dotted, teal!70!black, thick, domain=0.1:0.5, samples=30]
    plot (\x, {-\x*\x + \x});
    \node[teal!85!black, font=\scriptsize, above right] at (0.7, 0.2) {$\beta = -\alpha^2 + \alpha$};

    % --- 4. 境界条件の補助直線 beta = 1/2*alpha (破線) ---
    \draw[dashed, gray!60, thick] (0, 0) -- (0.9, 0.45)
    node[above right=-1pt, font=\scriptsize, black] {$\beta = \frac{1}{2}\alpha$};

    % --- 5. 座標軸 ---
    \draw[->, thick, gray!80!black] (-0.1, 0) -- (1.25, 0) node[right] {$\alpha$};
    \draw[->, thick, gray!80!black] (0, -0.08) -- (0, 0.48) node[above] {$\beta$};
    \node[below left=2pt] at (O) {$O$};

    % --- 6. 座標投影破線と目盛り ---
    % alpha = 1/2, beta = 1/4 の投影
    \draw[dashed, gray!50] (0.5, 0) -- (P_joint) -- (0, 0.25);
    \node[below=2pt, font=\scriptsize] at (0.5, 0) {$\frac{1}{2}$};
    \node[left=2pt, font=\scriptsize]  at (0, 0.25) {$\frac{1}{4}$};

    % alpha = 1 の目盛り
    \node[below=2pt, font=\scriptsize] at (1.0, 0) {$1$};

    % --- 7. 特徴点・除外点のプロット ---
    % 接続点 (1/2, 1/4) (境界に含まれる実点)
    \fill[blue!90!black] (P_joint) circle (0.6pt);
    \node[above=2pt, font=\scriptsize, blue!90!black] at (P_joint) {$\left(\frac{1}{2}, \frac{1}{4}\right)$};

    % 除外点 (alpha > 0, beta > 0 より軸上の境界点は含まない: 白抜き丸)
    \draw[thick, fill=white] (O) circle (0.7pt);
    \draw[thick, fill=white] (P_beta14) circle (0.7pt);
    \draw[thick, fill=white] (P_alpha1) circle (0.7pt);

  \end{tikzpicture}
  \caption{\cref{eq:6}の条件を満たす $(\alpha,\beta)$ の存在領域}
  \label{fig:alpha_beta_region_part1}
\end{figure}

これは\cref{eq:6}をまとめ直せば
\begin{align}
  \begin{cases}
    \beta \le \dfrac{1}{4} & \left(0 < \alpha \le \dfrac{1}{2}\right) \\
    \beta \le -\alpha^2+\alpha & \left(\dfrac{1}{2} \le \alpha < 1\right)
  \end{cases} \label{eq:7}
\end{align}
である．

\paragraph{(2)}

楕円$C$の面積は $T = \sqrt{\alpha\beta} \pi$ だから，$\alpha\beta$が最大の時に$T$も最大である．以下$\alpha$を固定して考える．\cref{eq:7}から$\alpha$が$\dfrac{1}{2}$より大きいかどうかで場合分けして考える．

\subparagraph{$1^\circ \ 0 \le \alpha \le \dfrac{1}{2}$}

この時は$0 < \beta \le \dfrac{1}{4}$ だから，$(\alpha, \beta) = \left(\dfrac{1}{2}, \dfrac{1}{4}\right)$ で最大値
\begin{align}
  \max \alpha\beta = \dfrac{1}{8} \label{eq:8}
\end{align}
をとる．

\subparagraph{$2^\circ \ \dfrac{1}{2} \le \alpha < 1$}

$\beta = -\alpha^2 + \alpha$ で $\alpha\beta$は最大値
\begin{align}
  \max \alpha\beta = \alpha(-\alpha^2 + \alpha) = -\alpha^3+\alpha^2
\end{align}
をとる．この右辺を$g(\alpha)$とおいて次に$\alpha$を動かす．
一階微分は
\begin{align}
  f'(\alpha) &= -3\alpha^2 + 2\alpha
\end{align}
だから$f(\alpha)$の増減表は以下のようになる．

\begin{table}[htb]
  \centering
  \caption{$f(\alpha)$の増減表}
  \begin{tabular}{|c|c|c|c|c|c|}
    \hline
    $\alpha$ & $\dfrac{1}{2}$ & $\cdots$ & $\dfrac{2}{3}$ & $\cdots$ & $1$ \\
    \hline
    $f'$ & & $+$ & $0$ & $-$ & \\
    \hline
    $f$ & & $\nearrow$ & & $\searrow$ & \\
    \hline
  \end{tabular}
\end{table}

したがって，$\alpha=\dfrac{2}{3}$の時に
\begin{align}
  \max \alpha\beta = f(2/3) = \dfrac{4}{27} \ (> \dfrac{1}{8})
\end{align}
をとる．これは$0 \le \alpha \le \dfrac{1}{2}$の時の最大値\cref{eq:8}よりも大きい．
\casesend

以上 $1^\circ, 2^\circ$ から $\max \alpha\beta = \dfrac{4}{27}$ だから，
求める$T$の最大値は
\begin{align}
  \max T
  = \pi\sqrt{\max \alpha\beta}
  = \dfrac{2}{9} \sqrt{3} \pi
\end{align}
である．

\end{document}
